\scanpage{53}
\deleted{Le \inserted{\deleted{\unsure{groupe algébrique}}} groupe engendré par les $\varphi(M(n))$ est donc défini,
c'est un ss-grpe algébrique lisse de $G$. Donc}

Supposons maintenant que $G$ soit un schéma en
groupes Néronien sur un trait $S$ de corps des
fonctions $K$, soit $M$ un grpe p-divisible, et
$\varphi_K:M_K\unsure{-|\to}G_K$. Je dis que $\varphi_K$ se prolonge en
$\varphi:M\to G$ (évidemment unique).

Clair si $p\ne\mathrm{car}.$ résiduelle.

Supposons donc $p=\mathrm{car}.\,k$. Alors, ou bien $\mathrm{car}.\,K=p$,
ou bien $\mathrm{car}.\,K=0$ (donc $K$ parfait); en tous cas,
$G_K$ n'a pas de ss-groupe \unsure{isom. à|isom. :} $G_a$, puisque
$G_K$ admet un modèle de Néron. Donc on
peut appliquer le lemme précédent, \deleted{Il existe} d'où
un $H_K\subset G_K$ par lequel se factorise $\varphi$.
Comme $H_K$ est ext. d'une VA par un tore,
$H_K$ admet un modèle de Néron $H$, et $H_K\overset{i_K}{\hookrightarrow}G_K$
est \unsure{induit|induite} par $H\xrightarrow{i}K$. Le th. de Raynaud
pour $H$ \unsure{implique|indique} que $\psi_K:M_K\to H_K$ se
prolonge en $\psi:M\longrightarrow H$. On gagne !
% archive sheet number: 52
